\chapter{Programacion Dinamica (dp)}

\section{De una función recursiva por partes a código con DP}

En muchos problemas, partimos de una definición de función recursiva por partes. El objetivo es transformarla en un algoritmo eficiente usando programación dinámica. A continuación se muestra un ejemplo general del proceso.

\subsection*{Ejemplo: función definida por partes}

Supongamos que tenemos la siguiente definición matemática:

\[
f(i) =
\begin{cases}
0 & \text{si } i = 0 \\
f(i - 1) + |a[i] - a[i-1]| & \text{si } i = 1 \\
\min(f(i - 1) + |a[i] - a[i-1]|,\ f(i - 2) + |a[i] - a[i-2]|) & \text{si } i \ge 2
\end{cases}
\]

Queremos implementarla eficientemente para valores grandes de \( i \).

\subsection*{Paso 1: Identificar el estado}

El subproblema está determinado solo por el índice \( i \). Por lo tanto, usamos un arreglo \texttt{dp[i]} donde \texttt{dp[i]} almacena el valor de \( f(i) \).

\subsection*{Paso 2: Código en C++ con memoización}

\begin{lstlisting}[language=C++, caption=Función recursiva con memoización]
const int MAXN = 1e5+10;
int dp[MAXN];
int a[MAXN];

int f(int i) {
if (dp[i] != -1) return dp[i];

if (i == 0) return dp[i] = 0;
if (i == 1) return dp[i] = f(0) + abs(a[1] - a[0]);

return dp[i] = min(
f(i - 1) + abs(a[i] - a[i - 1]),
f(i - 2) + abs(a[i] - a[i - 2])
);
}
\end{lstlisting}

\textbf{Inicialización:}

\begin{lstlisting}[language=C++]
fill(dp, dp + n, -1);
\end{lstlisting}

\textbf{Llamada inicial:}

\begin{lstlisting}[language=C++]
cout << f(n - 1) << '\n';
\end{lstlisting}

\subsection*{Paso 3: Versión iterativa (tabulación)}

\begin{lstlisting}[language=C++, caption=Versión iterativa (bottom-up)]
dp[0] = 0;
dp[1] = abs(a[1] - a[0]);

for (int i = 2; i < n; ++i) {
dp[i] = min(
dp[i - 1] + abs(a[i] - a[i - 1]),
dp[i - 2] + abs(a[i] - a[i - 2])
);
}
\end{lstlisting}

\subsection*{Resumen}

\begin{itemize}
\item La clave está en identificar bien los \textbf{casos base}.
\item Luego traducir cada parte de la definición en código con condiciones explícitas.
\item Asegurate de guardar los resultados para no recalcular.
\end{itemize}


\section{Ejemplos clásicos}

\subsection{A - Frog 1}

Problema: dado un arreglo de alturas, calcular el costo mínimo para saltar desde el índice 0 al último, con saltos de 1 o 2 posiciones.

\begin{lstlisting}[language=C++, caption=Frog1: DP mínima energía]
ll dp[MAXN];
ll a[MAXN];
ll f(int i) {
if (i == 0) return 0;
if (dp[i] != -1) return dp[i];
ll res = f(i-1) + abs(a[i] - a[i-1]);
if (i > 1) res = min(res, f(i-2) + abs(a[i] - a[i-2]));
return dp[i] = res;
}
\end{lstlisting}

Estados: posición actual.  
Transiciones: desde `i-1` o `i-2`.  
Caso base: `dp[0] = 0`.

\subsection{C - Vacation}

Problema: en cada día, podés hacer una de 3 actividades con recompensa. No podés repetir la misma actividad dos días seguidos.

\begin{lstlisting}[language=C++, caption=Vacation: DP sobre días y última actividad]
ll dp[MAXN][3];
ll a[MAXN][3];

ll f(int i, int last) {
if (i == n) return 0;
if (dp[i][last] != -1) return dp[i][last];
ll res = 0;
fore(j, 0, 3) {
if (j != last) {
res = max(res, a[i][j] + f(i+1, j));
}
}
return dp[i][last] = res;
}
\end{lstlisting}

Estado: día actual y actividad anterior.  
Transiciones: elegir una actividad distinta a la del día anterior.  
Caso base: cuando `i == n`.

\subsection{D - Knapsack 0/1}

Problema clásico: elegir un subconjunto de objetos sin pasar un peso máximo.

\begin{lstlisting}[language=C++, caption=Mochila 0/1: DP por peso]
ll dp[MAXW];
int w[MAXN], v[MAXN];

void solve() {
for (int i = 0; i < n; ++i) {
for (int j = W; j >= w[i]; --j) {
dp[j] = max(dp[j], dp[j - w[i]] + v[i]);
}
}
}
\end{lstlisting}

Estado: peso total actual.  
Transiciones: elegir o no cada objeto.  
Usa optimización de espacio (1D).

\subsection{G - Longest Path in DAG}

Problema: encontrar el camino más largo en un grafo dirigido acíclico.

\begin{lstlisting}[language=C++, caption=Camino más largo en DAG con DP + DFS]
vector<int> G[MAXN];
int dp[MAXN];

int dfs(int u) {
if (dp[u] != -1) return dp[u];
int res = 0;
for (int v : G[u]) {
res = max(res, 1 + dfs(v));
}
return dp[u] = res;
}
\end{lstlisting}

Estado: nodo actual.  
Transiciones: recorrer vecinos.  
Requiere topological sort o memoización con DFS.

---

\section{Resumen}

\begin{itemize}
\item Pensar en el subproblema mínimo que necesitás.
\item Formular la recurrencia primero.
\item Escribir la versión recursiva con memoización es buena forma de empezar.
\item Luego podés optimizar a iterativa si hace falta.
\end{itemize}